%===========================================
% ARQUIVO DE TESTE / TEMPLATE MODULAR MEMOIR
%===========================================
\documentclass[12pt,a4paper,oneside]{memoir}
\usepackage[brazilian]{babel}
\usepackage{lmodern}
\usepackage{tikz}
\usetikzlibrary{calc}
\usepackage{hyperref}
\usepackage{amssymb}  % Símbolos 
\usepackage{amsmath}  % Matemática
\usepackage{lipsum}
\usepackage{fontspec}
\usepackage{textcase}
\usepackage{etoolbox} 
\setmainfont{CMU Serif}


\usepackage[style=abnt,backend=biber]{biblatex}
\addbibresource{Arquivos/Referencias.bib}
\usepackage{ragged2e}
\AtBeginBibliography{\justifying}
\urlstyle{same}

\setlrmarginsandblock{3cm}{2cm}{*}
\setulmarginsandblock{3cm}{2cm}{*}
\checkandfixthelayout
\OnehalfSpacing
\setlength{\parindent}{1.25cm}







% ==========================================================
% sumário mostra ate subseção
% ==========================================================
\settocdepth{subsection}
\setsecnumdepth{subsection}
% ativar pontilhado nos capítulos no sumário
\renewcommand{\cftchapterdotsep}{\cftdotsep}
% ==========================================================

% ---------------------------------------------
% É controle espaço horizontal título sumário
% ---------------------------------------------
\cftsetindents{part}{0pt}{2cm}
\cftsetindents{chapter}{0pt}{2cm}
\cftsetindents{section}{0pt}{2cm}
\cftsetindents{subsection}{0pt}{2cm}
% ---------------------------------------------
 



% ==========================================================
% Defini capítulos para maiúsculas no sumario
% ==========================================================
\makeatletter
% Salva o comando original que desenha a linha do sumário
\let\oldcontentsline\contentsline

% Redefine para interceptar e modificar o texto
\renewcommand\contentsline[4]{%
  \ifboolexpr{
    test {\ifstrequal{#1}{chapter}} or 
    test {\ifstrequal{#1}{part}} or 
    test {\ifstrequal{#1}{appendix}} % <--- O SEGREDO: Pega também se for classificado como apêndice
  }
    {\oldcontentsline{#1}{\MakeTextUppercase{#2}}{#3}{#4}} % Se for um dos 3, Caixa Alta
    {\oldcontentsline{#1}{#2}{#3}{#4}} % Se for seção ou outro, mantém normal
}
\makeatother
% ==========================================================


% ==========================================================
% sumário
% ==========================================================

\makechapterstyle{SumarioStyle}{
  \setlength{\beforechapskip}{-73pt}  % ajuste fino da altura
  \setlength{\afterchapskip}{1\onelineskip}

  \renewcommand{\printchaptername}{}
  \renewcommand{\chapternamenum}{}
  \renewcommand{\printchapternum}{}

  \renewcommand{\printchaptertitle}[1]{%
    \vspace*{2\onelineskip}
    \begin{center}
      \Huge\bfseries ##1\par
    \end{center}%
  }
}
% ==========================================================





\newcommand{\Sumario}{
   \clearpage
   \pagestyle{empty}
   \aliaspagestyle{chapter}{empty}
   \chapterstyle{SumarioStyle}  % só para o sumário
   \tableofcontents*
   \cleardoublepage
   \aliaspagestyle{chapter}{plain}
   \chapterstyle{CapitulosNumerado}  % volta ao normal
}











%===========================================
% ESTILOS DE CAPÍTULO (MODULARES)
%===========================================
\makechapterstyle{CapitulosNumerado}{%
    \setlength{\beforechapskip}{0.46cm} % Descer estilo de capítulo
    \renewcommand{\printchaptername}{}
    \renewcommand{\chapternamenum}{}
    \renewcommand{\afterchapternum}{}
    \setlength{\midchapskip}{1.3cm} 
    \renewcommand{\afterchapternum}{\par\nobreak\vskip \midchapskip}
    \renewcommand{\printchapternum}{%
        \noindent
        \begin{tikzpicture}[remember picture, overlay]
            \def\numsize{61} % Tamanho do número          
            \def\barwidth{0.8pt} % Espessura da linha      
            \def\leftoffset{2.3cm} % Reta horizontal esquerda
            
            \node[inner sep=0pt, anchor=south west, opacity=0] (num) at (\leftoffset, 0)
             {\fontsize{\numsize}{\numsize}\selectfont\rmfamily\thechapter};
             
            \node[inner sep=0pt, anchor=south west] at (\leftoffset, -1pt)
             {\fontsize{\numsize}{\numsize}\selectfont\rmfamily\thechapter};
            
            \node[inner sep=0pt, anchor=south west, font=\large\rmfamily] (chaptext)
                 at (0pt, -4pt) {\chaptername};
            
            % --- ALTURA FIXA PARA O RETÂNGULO ---
            \def\boxheight{41pt}   % ajuste aqui até ficar perfeito
            
            % --- PONTO ESQUERDO DO L ---
            \coordinate (Lbase) at (0.5*\barwidth,15pt);         % base fixa
            \coordinate (Ltop)  at (0.5*\barwidth,\boxheight);   % topo fixo
            
            % --- DESENHA O L FIXO ---
            % --- AJUSTE: distância horizontal entre o L e o número ---
            \def\Lgap{0.2cm} % ajuste aqui até ficar visualmente perfeito
            
            % --- DESENHO DO L FIXO ---
            \draw[line width=\barwidth]
                (Lbase) -- (Ltop)
                -- ++(\leftoffset - 0.5*\barwidth - \Lgap,0);

            % --- PONTOS DO RETÂNGULO FIXO ---
            \coordinate (NumRightBase) at ($(num.south east)+(0.2cm,0)$);
            \coordinate (NumRightTop)  at ($(NumRightBase)+(0,\boxheight)$);
            
            \coordinate (RightTop)     at (\linewidth,\boxheight);
            \coordinate (RightBase)    at (\linewidth,0);
            
            % --- CAIXA FIXA ---
            \draw[line width=\barwidth]
                (NumRightTop) -- (RightTop)
                              -- (RightBase)
                              -- (NumRightBase);
        \end{tikzpicture}%
    }
    \renewcommand{\printchaptertitle}[1]{%
        \raggedright\huge\bfseries ##1
    }
}




%===========================================

\makechapterstyle{estiloNaoNumerado}{%
  \setlength{\beforechapskip}{-73pt}%
  \setlength{\afterchapskip}{2\onelineskip}%
  \renewcommand{\printchaptername}{}%
  \renewcommand{\chapternamenum}{}%
  \renewcommand{\printchapternum}{}%
  \renewcommand{\printchaptertitle}[1]{%
    \vspace*{2\onelineskip}%
    \begin{center}
      \Huge\bfseries ##1\par
      \vspace{-1.5\onelineskip}%
      \rule{\textwidth}{0.4pt}%
    \end{center}%
  }%
}
%===========================================




%===========================================
\def\defaultchapterstyle{CapitulosNumerado}

\makeatletter
\newcommand{\capituloCentralizado}{%
  \@ifstar{\capituloCentralizado@star}{\capituloCentralizado@nostar}}
  
% VERSÃO NORMAL – entra no sumário e segue o comportamento atual
\newcommand{\capituloCentralizado@nostar}[1]{%
  \chapterstyle{estiloNaoNumerado}%
  \chapter*{#1}%
  \markboth{#1}{}%
  \addcontentsline{toc}{chapter}{\protect\numberline{}#1}
  \chapterstyle{CapitulosNumerado}%
}

% VERSÃO COM * – sem número de página e FORA do sumário
\newcommand{\capituloCentralizado@star}[1]{%
  \chapterstyle{estiloNaoNumerado}%
  \chapter*{#1}%
  \thispagestyle{empty}%              % sem número de página
  \markboth{#1}{}%
  \chapterstyle{CapitulosNumerado}%
}
\makeatother

%===========================================







%===========================================

% LISTA DE FIGURAS (pré-textual, sem número de página, fora do sumário)
\newcommand{\listaDeFiguras}{%
  \clearpage
  \chapterstyle{estiloNaoNumerado}%
  \listoffigures*% usa o estilo acima e NÃO entra no sumário
  \thispagestyle{empty}% tira o número da página desta folha
  \chapterstyle{CapitulosNumerado}%
}
%===========================================

% LISTA DE TABELAS (mesma ideia)
\newcommand{\listaDeTabelas}{%
  \clearpage
  \chapterstyle{estiloNaoNumerado}%
  \listoftables*%
  \thispagestyle{empty}%
  \chapterstyle{CapitulosNumerado}%
}
%===========================================



%===========================================
\newcommand{\preTexto}[1]{%
  \clearpage
  \chapterstyle{estiloNaoNumerado}%
  \chapter*{#1}%
  \thispagestyle{empty}%
  \chapterstyle{CapitulosNumerado}%
}

% Lista de abreviaturas e siglas
\newcommand{\listadesiglasname}{Lista de abreviaturas e siglas}
\newenvironment{siglas}{%
  \preTexto{\listadesiglasname}
  \begin{symbols}
}{%
  \end{symbols}
  \cleardoublepage
}

% Lista de simbolos
\newcommand{\listadesimbolosname}{Lista de símbolos}
\newenvironment{simbolos}{%
  \preTexto{\listadesimbolosname}
  \begin{symbols}
}{%
  \end{symbols}
  \cleardoublepage
}
%===========================================



% Exibir nome (Figuras, Tabelas) na lista 
\renewcommand{\cftfigurename}{Figura\space}
\renewcommand{\cfttablename}{Tabela\space}
% Exibir "-" na lista
\renewcommand{\cftfigureaftersnum}{\space\space--}
\renewcommand{\cfttableaftersnum}{\space\space--}

\cftsetindents{figure}{0em}{3em}
\cftsetindents{table}{0em}{3em}










%===========================================
% CABEÇALHOS E NUMERAÇÃO
%===========================================
\makepagestyle{meuestilo}
\makeevenhead{meuestilo}{\footnotesize\textit{\leftmark}}{}{\footnotesize\thepage}
\makeoddhead {meuestilo}{\footnotesize\textit{\leftmark}}{}{\footnotesize\thepage}
\makeheadrule{meuestilo}{\textwidth}{0.3pt}

\makepagestyle{plain}
\makeoddhead{plain}{}{}{\footnotesize\thepage}
\makeevenhead{plain}{}{}{\footnotesize\thepage}
\makeheadrule{plain}{\textwidth}{0pt}

\makeatletter
\renewcommand{\chaptermark}[1]{%
  \markboth{\thechapter.\ #1}{}%
}
\makeatother

%===========================================
% FRONT / MAIN / BACK MATTER
%===========================================
\makeatletter
\renewcommand{\frontmatter}{%
  \cleardoublepage
  \pagenumbering{arabic}%
  \pagestyle{empty}%
  \chapterstyle{estiloNaoNumerado}%
}
\renewcommand{\mainmatter}{%
  \cleardoublepage
  \pagestyle{meuestilo}%
  \chapterstyle{\defaultchapterstyle}%
}
\renewcommand{\backmatter}{%
  \cleardoublepage
  \pagestyle{meuestilo}%
  \chapterstyle{\defaultchapterstyle}%
}
\makeatother




\newcommand{\paginaDivisoria}[1]{%
  \clearpage
  \thispagestyle{empty}
  \phantomsection % Cria âncora para o link 
  \vspace*{8cm}
  \begin{center}
     \textbf{\Huge #1}
  \end{center}
  \addcontentsline{toc}{part}{\protect\numberline{}\large\textbf{#1}}
  \clearpage
}




\newcommand{\inicioApendices}{%
  \appendix
  \setcounter{chapter}{0}%
  \renewcommand{\chaptername}{Apêndice}%
  \renewcommand{\appendixname}{APÊNDICE}%
  \renewcommand{\thechapter}{\Alph{chapter}}%
  \makeatletter
  \renewcommand{\chaptermark}[1]{%
    \markboth{APÊNDICE~\thechapter.\ ##1}{}%
  }%
  \makeatother
}

\newcommand{\inicioAnexos}{%
  \cleardoublepage
  \setcounter{chapter}{0}%
  \renewcommand{\chaptername}{Anexo}%
  \renewcommand{\thechapter}{\Alph{chapter}}%
  \makeatletter
  \renewcommand{\chaptermark}[1]{%
    \markboth{ANEXO~\thechapter.\ ##1}{}%
  }%
  \makeatother
}



\cftsetindents{chapter}{0pt}{2cm}
\cftsetindents{section}{0pt}{2cm}
\cftsetindents{subsection}{0pt}{2cm}



%===========================================
% FOLHA / CONTRA
%===========================================
\newcommand{\folha}{%
    \thispagestyle{empty}%
    \vspace*{10em}%
    \begin{center}
    {\bfseries\scshape\Huge Folha de Rosto \par}
    \end{center}%
    \clearpage
}

\newcommand{\contra}{%
  \thispagestyle{empty}%
  \vspace*{8cm}%
  \begin{center}
    {\bfseries\scshape\Huge Contra Capa \par}
  \end{center}%
  \clearpage
}














%===========================================
% DOCUMENTO
%===========================================
\begin{document}
\nocite{*}
\frontmatter

\folha
\contra

\capituloCentralizado*{Agradecimentos}
\lipsum[1]

\capituloCentralizado*{Resumo}
\lipsum[1]

\listaDeFiguras
\listaDeTabelas


\begin{siglas}
  \item[CPU] Unidade Central de Processamento
  \item[RAM] Memória de Acesso Aleatório
  \item[SSD] Unidade de Estado Sólido (Solid State Drive)
  \item[OS] Sistema Operacional (Operating System)
  \item[HTTP] Protocolo de Transferência de Hipertexto
  \item[SQL] Linguagem de Consulta Estruturada (Structured Query Language)
  \item[API] Interface de Programação de Aplicações (Application Programming Interface)
  \item[ISP] Provedor de Serviços de Internet (Internet Service Provider)
\end{siglas}

\begin{simbolos}
\item[$\alpha$] Alfa (Ângulo, Coeficiente de Temperatura)
\item[$\beta$] Beta (Ângulo, Ganho de Corrente do Transistor)
\item[$\sum$] Somatório (Usado em Séries Matemáticas e Estatística)
\item[$\int$] Integral (Usado em Cálculo para Áreas ou Acúmulo)
\item[$\nabla$] Nabla (Usado em Cálculo Vetorial, como Gradiente)
\item[$\Omega$] Ohm (Unidade de Resistência Elétrica; Símbolo para Computação Quântica)
\item[$\approx$] Aproximadamente Igual a (Usado para Estimativas)
\item[$\implies$] Implica (Usado em Lógica Matemática)
\end{simbolos}




\Sumario

\mainmatter

\capituloCentralizado{Introdução}
\lipsum[1]

\part{Preparação do Relatório}


\chapter{Desenvolvimento}
\section{Coleta de Dados}
\section{Desenvolvimento}
\subsection{Desenvolvimento}
\lipsum[1-5]
\subsubsection{Coleta de Dados}

\begin{figure}[ht]
    \centering
    \begin{tikzpicture}
        \draw[rounded corners=8pt, thick] (0,0) rectangle (5,3);
        \node at (2.5,1.5) {Figura de Teste};
    \end{tikzpicture}
    \caption{Figura de teste estilizada}
    \label{fig:teste-estilosa}
\end{figure}
\begin{figure}[ht]
    \centering
    \begin{tikzpicture}
        \draw[thick, fill=blue!20] (2.5, 1.5) circle (1.2);
        \draw[ultra thick, ->, >=stealth, red!70!black] (2.5, 2.7) -- (2.5, 3.8);
        \node at (2.5,1.5) {Processo};
        \node at (2.5,4.1) {Saída};
    \end{tikzpicture}
    \caption{Fluxo circular simples.}
    \label{fig:fluxo-simples}
\end{figure}
\begin{figure}[ht]
    \centering
    \begin{tikzpicture}
        \draw[rounded corners=5pt, very thick, fill=green!15] (0, 0) rectangle (2, 2);
        \node at (1, 1) {Entrada};
        \draw[rounded corners=5pt, very thick, fill=orange!15] (3, 0) rectangle (5, 2);
        \node at (4, 1) {Resultado};
        \draw[->, very thick, black!70] (2, 1) -- (3, 1);
        \node at (2.5, 1.3) {Função};
    \end{tikzpicture}
    \caption{Diagrama de blocos interligados.}
    \label{fig:blocos-interligados}
\end{figure}
\begin{figure}[ht]
    \centering
    \begin{tikzpicture}[scale=0.8]
        \draw[->, thick] (0,0) -- (5,0) node[right] {Eixo $x$};
        \draw[->, thick] (0,0) -- (0,4) node[above] {Eixo $y$};
        \fill (0,0) circle (2pt) node[below left] {O};
        \draw[ultra thick, blue!70] (0,0) parabola (4,3);
        \fill[red] (2, 0.75) circle (3pt) node[above right] {$(x_1, y_1)$};
    \end{tikzpicture}
    \caption{Gráfico de coordenadas cartesianas.}
    \label{fig:grafico-cartesianas}
\end{figure}



\part{Resultado do Relatório}



\chapter{Conclusão}
\section{Donec pellentesque}
\lipsum[55]
\section{Donec pellentesque}
\lipsum[101]
\begin{table}[ht]
    \centering
    \caption{Tabela estilosa 5×5}
    \label{tab:estilosa}
    \begin{tabular}{ccccc}
        \toprule
        Col 1 & Col 2 & Col 3 & Col 4 & Col 5 \\
        \midrule
        A1 & B1 & C1 & D1 & E1 \\
        A2 & B2 & C2 & D2 & E2 \\
        A3 & B3 & C3 & D3 & E3 \\
        A4 & B4 & C4 & D4 & E4 \\
        A5 & B5 & C5 & D5 & E5 \\
        \bottomrule
    \end{tabular}
\end{table}
\begin{table}[ht]
    \centering
    \caption{Comparativo de Tecnologias de Rede}
    \label{tab:comp-redes}
    \begin{tabular}{lcc}
        \toprule
        Tecnologia & Alcance Típico & Velocidade Máxima (Teórica) \\
        \midrule
        Ethernet (Gigabit) & Curto & 1 Gbps \\
        Wi-Fi 6 (802.11ax) & Médio & 9.6 Gbps \\
        Fibra Óptica & Longo & > 100 Gbps \\
        Bluetooth & Muito Curto & 3 Mbps \\
        \bottomrule
    \end{tabular}
\end{table}

\begin{table}[ht]
    \centering
    \caption{Componentes FV com Quebra de Linha e Largura Fixa}
    \label{tab:comp-fv-fixo}
    \begin{tabular}{|p{3cm}|p{3cm}|c|c|}
        \toprule
        \textbf{Componente} & \textbf{Função} & \textbf{Tipo} & \textbf{Medida} \\
        \midrule
        Painel PV & Geração & CC (DC) & Wp \\
        \midrule
        Inversor & Conversão & Entrada & VA \\
        \midrule
        Controlador & Bateria & CC (DC) & A \\
        \midrule
        Bateria  & Armazenamento & CC (DC) & Ah \\
        \bottomrule
    \end{tabular}
\end{table}


\begin{table}[ht]
    \centering
    \caption{Classificação Básica de Conjuntos Numéricos}
    \label{tab:conjuntos-mat}
    \begin{tabular}{cl}
        \toprule
        Símbolo & Descrição do Conjunto \\
        \midrule
        $\mathbb{N}$ & Naturais (Contagem: 1, 2, 3, ...) \\
        $\mathbb{Z}$ & Inteiros (Naturais + Negativos + Zero) \\
        $\mathbb{Q}$ & Racionais (Números que podem ser escritos como fração) \\
        $\mathbb{R}$ & Reais (Todos os números racionais e irracionais) \\
        $\mathbb{C}$ & Complexos (Reais + Imaginários) \\
        \bottomrule
    \end{tabular}
\end{table}


\section{Donec pellentesque}
\lipsum[101]
\section{Donec pellentesque}
\lipsum[101]


\backmatter

\capituloCentralizado{Referências}
\printbibliography[heading=none]

\paginaDivisoria{APÊNDICES}
\inicioApendices

\chapter{Título do Apêndice}
\section{Conclusão}
\lipsum[1]
\section{Donec pellentesque}
\lipsum[101]


\chapter{Título do Apêndice}
\section{Conclusão}
\lipsum[1]
\section{Donec pellentesque}
\lipsum[101]

\paginaDivisoria{ANEXOS}
\inicioAnexos

\chapter{Título do Anexo}
\section{Conclusão}
\lipsum[1]
\section{Donec pellentesque}
\lipsum[101]



\chapter{Título do Anexo}
\section{Conclusão}
\lipsum[1]
\section{Donec pellentesque}
\lipsum[101]
\end{document}
